\chapter{API \& integrations}
\label{ch:api}

mybibli ships a small JSON HTTP API at \texttt{/api/v1/*} that lets
external scripts and AI tools read your catalog and update the
classification of titles. The API runs on the same web server as the
human-facing UI — there is no separate daemon, no extra port to open.

This chapter walks through:

\begin{itemize}
  \item creating an API key in the admin panel,
  \item the two scopes (\emph{read only} vs. \emph{read + write}),
  \item the endpoints exposed under \texttt{/api/v1/*},
  \item three end-to-end examples (\texttt{curl}, Python, an AI
        classifier loop).
\end{itemize}

\section{Why an HTTP API?}

The everyday use case is offloading classification work to an
external tool — for instance, asking an LLM to suggest a Dewey
code, a genre, or a clean subtitle for every newly-added English
book. The LLM walks the catalog over \texttt{GET /api/v1/titles},
fetches the detail of each via \texttt{GET /api/v1/titles/\{id\}},
computes a suggestion, and writes it back with
\texttt{PATCH /api/v1/titles/\{id\}}. Every PATCH is audited so you
can roll back changes if a model produced garbage.

The API is \emph{not} meant for high-throughput batch ingest — for
that, use the SQL database directly during a maintenance window.
It \emph{is} meant for the kind of slow, careful, opinionated
write that fits inside an LLM agent loop.

\section{Creating an API key}
\label{sec:api-keys-create}

API keys are managed under \textbf{Administration \(\rightarrow\)
API keys}. Only Admin users see this tab.

\begin{enumerate}
  \item Click \textbf{Administration} in the navigation bar.
  \item Open the \textbf{API keys} tab.
  \item Fill in a \textbf{label} — a name only you see, e.g.
        \emph{``AI classifier''} or \emph{``Backup script''}. The
        label is for the audit log; the key plaintext itself is
        random and unrelated to it.
  \item Choose a \textbf{scope}:
        \begin{itemize}
          \item \textbf{Read only} — \texttt{GET} endpoints only.
                Safe for read-only dashboards, mirroring tools,
                or curiosity-driven scripts.
          \item \textbf{Read + write} — additionally allows
                \texttt{PATCH /api/v1/titles/\{id\}} for a
                narrow allow-list of fields (genre, dewey,
                description, subtitle). Use this for LLM agents
                or curation tools.
        \end{itemize}
  \item Click \textbf{Create key}. A modal appears showing the
        full plaintext exactly once — copy it now. Once you
        close the modal there is no recovery path; if you lose
        the key, revoke it and create a new one.
\end{enumerate}

The plaintext shape is
\texttt{mybibli\_<scope>\_<random40>}, where \texttt{<scope>} is
\texttt{ro} or \texttt{rw}. Only the argon2 hash and the first
12 characters of the plaintext (the prefix) are persisted; the
prefix is what mybibli uses to narrow candidates during
authentication. The hash itself cannot be reversed.

\section{Authenticating}

Every \texttt{/api/v1/*} request must carry one of:

\begin{itemize}
  \item \texttt{Authorization: Bearer <plaintext>} — preferred.
  \item \texttt{X-API-Key: <plaintext>} — convenience fallback.
\end{itemize}

A missing or invalid key returns \texttt{401 Unauthorized} with
a JSON envelope:

\begin{verbatim}
{ "error": "unauthorized",
  "reason": "missing_api_key" }
\end{verbatim}

A read-only key on a write endpoint returns \texttt{403
Forbidden} with the missing scope named:

\begin{verbatim}
{ "error": "forbidden",
  "reason": "insufficient_scope",
  "scope_required": "write" }
\end{verbatim}

CSRF protection is intentionally bypassed for \texttt{/api/*}
because bearer-token authentication does not ride on cookies and
the cross-site replay risk that CSRF defends against does not
apply.

\section{Endpoints}

\subsection{GET /api/v1/titles}

Paginated list. Supports \texttt{?page=N},
\texttt{?q=<search>}, \texttt{?genre\_id=N}, and
\texttt{?dewey\_prefix=8}.

Returns:

\begin{verbatim}
{ "items": [ { "id": 42, "title": "...",
               "subtitle": null,
               "language": "en",
               "media_type": "book",
               "isbn": "9780...",
               "publisher": "...",
               "publication_year": 2020,
               "genre_id": 7,
               "genre_name": "Fiction",
               "dewey_code": "813.54" },
             ... ],
  "page": 1, "total_pages": 12,
  "total_items": 287 }
\end{verbatim}

\subsection{GET /api/v1/titles/\{id\}}

Full detail, including contributors, volumes (with their
location), and series assignments. The
\texttt{version} field is the optimistic-locking handle —
keep it next to the data your script will edit.

\subsection{GET /api/v1/genres, /locations, /series}

Reference data. Use these when an LLM needs to know which genres
exist before it suggests a \texttt{genre\_id}.

\subsection{GET /api/v1/dewey/\{prefix\}}

Bucket lookup. Pass a partial Dewey prefix (e.g.
\texttt{8} for literature) and get back every title whose
\texttt{dewey\_code} starts with it. Useful for ``find all books
that already have a Dewey code in the 800s so I can compare''.

\subsection{PATCH /api/v1/titles/\{id\}}

\textbf{Write-scope only.} Updates the four whitelisted fields:

\begin{itemize}
  \item \texttt{subtitle} — nullable.
  \item \texttt{description} — nullable.
  \item \texttt{dewey\_code} — nullable; digits, dots, slashes only.
  \item \texttt{genre\_id} — must reference an existing,
        non-deleted genre.
\end{itemize}

The JSON body uses partial-update semantics:

\begin{itemize}
  \item Field \textbf{omitted} \(\rightarrow\) no change.
  \item Field set to \texttt{null} \(\rightarrow\) clear to
        NULL (for the three nullable fields above; \texttt{null}
        on \texttt{genre\_id} returns \texttt{400}).
  \item Field set to a value \(\rightarrow\) update.
\end{itemize}

\texttt{version} is required. On mismatch the response is
\texttt{409 Conflict} carrying both the expected and the
supplied version — refetch and retry.

A successful PATCH writes one row into \texttt{admin\_audit}
with action \texttt{api\_patch\_title}, attribution via the
issuing admin (\texttt{user\_id}) and the API key
(\texttt{details.via\_api\_key\_id} +
\texttt{details.via\_api\_key\_label}), and a per-field
\texttt{old}/\texttt{new} diff. A no-op PATCH (no field
actually changed) returns \texttt{200} but does NOT bump
\texttt{version} and does NOT write an audit row.

\section{Examples}

\subsection{\texttt{curl}}

\begin{verbatim}
KEY="mybibli_ro_..."

# List the first page of titles.
curl -H "Authorization: Bearer $KEY" \
     https://library.example.org/api/v1/titles

# Show one title in full.
curl -H "Authorization: Bearer $KEY" \
     https://library.example.org/api/v1/titles/42

# PATCH a Dewey code (needs a write-scope key).
KEY_WRITE="mybibli_rw_..."
curl -X PATCH \
     -H "Authorization: Bearer $KEY_WRITE" \
     -H "Content-Type: application/json" \
     -d '{"version": 3, "dewey_code": "813.54"}' \
     https://library.example.org/api/v1/titles/42
\end{verbatim}

\subsection{Python}

\begin{verbatim}
import os, requests

BASE = "https://library.example.org/api/v1"
KEY  = os.environ["MYBIBLI_API_KEY"]
H    = {"Authorization": f"Bearer {KEY}"}

# Walk every page.
page = 1
while True:
    r = requests.get(f"{BASE}/titles?page={page}",
                     headers=H, timeout=30)
    r.raise_for_status()
    data = r.json()
    for t in data["items"]:
        print(t["id"], t["title"], t.get("dewey_code"))
    if page >= data["total_pages"]:
        break
    page += 1
\end{verbatim}

\subsection{An LLM classifier loop}

The intended write-flow looks like this — pseudocode:

\begin{verbatim}
for t in walk_titles(KEY_READ):
    if t["dewey_code"]:
        continue                   # already classified
    detail = get(f"/titles/{t['id']}")
    suggestion = ask_llm(detail)   # returns {"dewey_code": ...}
    if not suggestion.get("dewey_code"):
        continue
    patch(
        f"/titles/{t['id']}",
        json = {
            "version":    detail["version"],
            "dewey_code": suggestion["dewey_code"],
        },
        key  = KEY_WRITE,
    )
\end{verbatim}

If two classifiers race on the same title, the loser gets a
\texttt{409} and re-fetches; no silent clobber. The audit log
records who (key) suggested what — handy for rolling back a
model that has a bad day.

\section{Revoking a key}

Revocation is one click in the admin panel. A revoked key
\emph{cannot} be restored; the audit history of past usage is
preserved alongside the now-defunct row. If you suspect a key has
leaked (committed to a public repo, e-mailed by mistake, found in a
log file), revoke immediately and create a new one.

\section{What the API does \emph{not} expose}

By design, the v1 API is read-mostly. It does not:

\begin{itemize}
  \item create or delete titles, volumes, contributors, borrowers,
        loans, or series;
  \item update any field outside the four-field PATCH allow-list;
  \item return passwords, password hashes, session tokens, or any
        other credentials;
  \item accept multipart uploads (no cover-image PUT).
\end{itemize}

These are intentional v1 limits — the use case that justified the
API is classification, and that fits inside the narrow allow-list.
Wider write-surface would change the security posture (CSRF
considerations, mass-mutation rate limits) and is deferred until a
concrete need surfaces.
